\section{Discussion}\label{sec:discussion}

\subsection{What was learned about mechanism}

The theoretical half of the paper reduced ``why do swimmers fade'' to a
sign-and-shape question, and the data answered the sign part cleanly. In
345 official races the profile is positively split, the reserve-coupled
mechanism M2 is sign-contradicted and remains the worst model even after
its coupling is fitted to the bound, and the even class is rejected on the
held-out set by a margin of 0.20 percentage points with an interval that
does not approach zero. Whatever makes late yards dearer than early ones in
a 200-yard freestyle, it is not the value of a full reserve. The positive
sign survives every credit in the registered band but the last; the
rejection of the even class holds for credits up to 2.6~s, which covers
the elite-anchored value and everything below it but not the field-pace
values above it. Within that qualification, it is the empirical result we
are most confident of.

The shape part was not answered, and the reason is instructive. M3 and M4
tie on held-out accuracy because they err in opposite directions on the
one statistic that was supposed to separate them: the observed fade is
front-loaded in kind, as M4 predicts, but only about two thirds as
front-loaded as the fitted M4 requires and far more than the fitted M3
allows. A mechanism between them --- a cost that rises with position
\emph{and} a ceiling that binds early --- would fit better, but adding it
would be a third one-parameter model chosen after the data were seen, and
we decline to do that here. The honest statement is that a single
one-parameter mechanism of either kind captures the mean profile to within
about half a percentage point per lap, that the residual is the finishing
kick, and that the fade's distribution between the early and late laps is
the next thing a model of this class must get right.

\subsection{The start credit is the problem}

The pilot showed the start-credit confound qualitatively; the registered
evaluation quantifies it. Across the literature-anchored band the fitted
$\beta_x$ moves by a factor of thirty and the held-out ranking passes
through three different winners. The fatigue mechanisms differ from each
other by tenths of a percentage point per lap; a tenth of a second of dive
credit moves the corrected first-lap share by about a tenth of a point.
This is not a statistical limitation to be cured by more races --- the
confidence intervals on the held-out RMSEs are already narrow relative to
the model differences at any fixed credit --- but a modelling limitation:
the free-swimming model is asked to explain a lap that is not free swimming,
and how much of that lap to credit to the dive is unknown to within a
second. The plan's start-uncertainty clause was written for exactly this
outcome, and we report it as the clause requires. Two remedies follow, in
order of value. First, measure the credit: 15~m split times are recorded at
some meets and published at almost none, and a single meet with them would
pin the credit for its field to a few tenths. Second, model it: a
pace-dependent term $c(v)=15/v-t_{15}$, with $t_{15}$ from the start
literature, replaces a constant that is wrong in a known direction for
slower swimmers and would remove the artefactual part of the ``faster
swimmers pace more evenly'' pattern before it is interpreted.

\subsection{The finishing kick}

Half the races in the dataset have a fourth lap faster than the third, and
the mean profile has its slowest lap third. No mechanism in the family can
produce this, and the reason is structural rather than parametric: every
cost multiplier we consider is monotone in position or reserve, and every
ceiling tightens monotonically, so the optimal lap times are monotone.
Three explanations are compatible with the data and none is tested by it.
A swimmer racing under uncertainty about their own reserve rationally holds
a margin until the finish is certain and then spends it; that is an
optimal-control problem with a stochastic state, not a deterministic one,
and it produces an end spurt. A swimmer racing a field rather than a clock
responds to the swimmers beside them, which the model excludes by design.
And a swimmer's underwater work off the last turn may simply be better,
because there is nothing to save it for. The first explanation is the most
interesting because it is a model, and it predicts something testable: the
size of the kick should fall as a swimmer's experience of the event grows
and their uncertainty about their state shrinks.

\subsection{Deviation and performance}

The pre-registered null on H1 was the expected outcome at this sample size
and it arrived. It should be read narrowly. The theory predicts that a 5\%
opening error costs 0.09~s in a 100~s race; the race-to-race noise in a
developing 16-year-old's 200 free is an order of magnitude larger, and the
expectation model available to us --- the swimmer's best earlier race in
this dataset --- is weaker still, as the three four-year-old personal bests
that manufactured a spurious curvature showed. A null under these
conditions does not say that pacing does not matter; it says that the
effect predicted by the model is below what this design can see, which the
plan said in advance. A design that could see it would use within-swimmer
comparisons across many races of one athlete, with seed times or
season-best trajectories as the expectation, and a deviation measured from
that swimmer's own fitted shape rather than the population's.

\subsection{What the theory contributes regardless}

The empirical section cannot yet say which fatigue mechanism is true, but
the structural results do not depend on it. Even pacing is the unique
optimum of every budget-constrained model with a velocity-only cost, so a
fade is always evidence of a mechanism and never a consequence of fitness;
the optimal shape is invariant to the engine and the hull, so no swimmer
has earned a faster opening lap by being stronger; and the penalty for
pacing error is quadratic and small near the optimum, so the difference
between a good pacing plan and the best one is hundredths of a second while
the difference between a good one and a bad one is tenths. The elite
profile of \citet{robertson2009analysis}, once the dive is credited, costs
0.01~s against the model's optimum; the mean profile of our first pilot
costs about 0.08~s; the ``very aggressive'' opening that coaches warn
against costs nearly a second. Which small error is cheapest depends on
the mechanism, but the ordering of large errors against small ones holds
for every mechanism considered and every parameter draw tried, and that is
the practical content of the paper.
